\documentclass[11pt]{article}
\usepackage[utf8]{inputenc}
\usepackage[T1]{fontenc}
\usepackage[a4paper,margin=2.5cm]{geometry}
\usepackage{amsmath,amssymb}
\usepackage{booktabs}
\usepackage{hyperref}
\usepackage{enumitem}
\hypersetup{colorlinks=true,linkcolor=blue,citecolor=blue,urlcolor=blue}

\title{\textbf{Brain-predictivity as the honest test:\\
the Cichy-92 benchmark for the G\"omb--Soma vision model}}
\author{Csaba Hajdu \\ \small Kato Laboratory \quad (working note, 2026-06-15)}
\date{}

\begin{document}
\maketitle

\begin{abstract}
\noindent
The G\"omb--Soma vision stack --- a signed-hypergraph object model built on an
adaptive quadtree, Forman--Ricci curvature, Hodge decomposition and
curvature-driven rewiring --- does \emph{not} win at small-scale image
classification: in a fair re-benchmark it trails even a plain MLP on MNIST and
Fashion-MNIST. That is a falsification of the \emph{accuracy} hypothesis, not of
the model's reason to exist. This note argues that the right test for a
structurally-grounded, brain-inspired vision model is not classification accuracy
but \emph{brain-predictivity}: how well its internal representations explain
measured cortical activity. We describe the Cichy~92-image MEG/fMRI benchmark
\cite{cichy2014}, the representational-similarity and Brain-Score methodology
\cite{kriegeskorte2008,schrimpf2018}, the parameter-matched protocol already
implemented in this repository, and the single missing piece --- the real fMRI
data --- needed to run it.
\end{abstract}

\section{Why this benchmark, and why now}
A vision model can be judged on two very different axes. The first is
\emph{task performance} (does it classify/detect well?). On this axis the
G\"omb--Soma family is, honestly, a negative result at the scales we can afford:
the 2026-05-28 re-benchmark (CNN, MLP, HSiKAN, NeocogHGNN, RicciStim on
MNIST/Fashion-MNIST, 15 epochs $\times$ 3 seeds) places RicciStim and the
hypergraph operators \emph{below} a vanilla MLP, and the differential-geometric
machinery does not recover the gap. The structural prior, on natural-image
classification at small scale, hurts.

The second axis is \emph{representational alignment with biological vision}: a
model is interesting to neuroscience if its layer activations linearly predict
the responses of visual cortex, \emph{independently} of whether it tops a
classification leaderboard. Many high-accuracy networks are poor cortical
predictors, and some lower-accuracy models align surprisingly well; accuracy and
brain-predictivity are related but distinct \cite{yamins2014,schrimpf2018}. Since
G\"omb--Soma is explicitly motivated by cortical structure --- coarse-to-fine
receptive fields, signed lateral interactions, curvature/flow on a retinotopic
graph --- the fair question is the second one. The Cichy-92 benchmark is the
standard, compact instrument for asking it.

\section{The Cichy 92-image dataset}
Cichy, Pantazis and Oliva \cite{cichy2014} presented 92 isolated object images
(human bodies, faces, animals, natural and artificial objects) to human
participants while recording \textbf{both} MEG (millisecond temporal resolution)
and fMRI (millimetre spatial resolution). The two modalities are fused through
\emph{representational similarity analysis} (RSA) \cite{kriegeskorte2008}: each
modality, region and time point yields a $92\times 92$ \emph{representational
dissimilarity matrix} (RDM) over the image pairs, and models are compared to
brains in this common RDM space rather than in raw activations. For the spatial
(fMRI) side the canonical regions of interest are early visual cortex
(V1) through higher object-selective cortex (e.g.\ IT/LOC), which makes the
dataset a natural probe for a \emph{hierarchical} model: a coarse-to-fine feature
stack should map onto the early-to-late cortical gradient.

The set is small (92 conditions), which is a feature: it is large enough for
stable RDM estimation with split-half noise ceilings, and small enough that a
parameter-light model and a parameter-matched baseline can be compared
end-to-end in seconds-to-minutes per seed.

\section{Scoring methodology (Brain-Score protocol)}
We follow the Brain-Score pipeline \cite{schrimpf2018}, implemented here without
the external \texttt{brainscore} package (no network dependency):
\begin{enumerate}[nosep]
  \item \textbf{Feature extraction.} For each of the 92 images, take a model's
        per-stage activations and bin them retinotopically (per quadtree depth
        for G\"omb--Soma; per residual block for the ResNet baseline).
  \item \textbf{Dimensionality reduction.} PLS regression compresses model
        features toward the neural target, guarding against the
        high-dimensional-features / few-conditions regime.
  \item \textbf{Per-voxel / per-ROI ridge mapping.} A cross-validated linear
        (ridge) map from model features to each ROI's responses; the score is the
        held-out predictivity (Pearson $r$ or $R^2$).
  \item \textbf{Noise-ceiling correction.} Split-half reliability of the neural
        data bounds the achievable score; reported predictivity is normalised by
        this ceiling so the number is interpretable as ``fraction of explainable
        variance captured''.
\end{enumerate}
The comparison is \emph{parameter-matched}: a $\sim$30k-parameter ResNet-tiny is
scored against the $\sim$30k-parameter G\"omb--Soma backbone on the identical
pipeline, so any difference reflects representational geometry rather than
capacity.

\section{The G\"omb--Soma cortical hypothesis}
The model offers three concrete, cortically-motivated features whose
brain-predictivity is testable:
\begin{itemize}[nosep]
  \item \textbf{Per-depth anchors as a V1$\to$V4 gradient.} The adaptive quadtree
        subdivides the image where salience/variance is high; shallow depths are
        coarse and uniform (V1-like), deep depths are fine and content-driven
        (V4/IT-like). Per-depth binned features are the model's ``cortical
        stages''.
  \item \textbf{Forman--Ricci curvature \cite{forman2003}} on the patch graph as a
        boundary/edge descriptor --- a candidate correlate of contour and
        figure-ground signals in early/mid visual areas.
  \item \textbf{Hodge eigenmodes} of the patch-graph Laplacian as global shape
        descriptors, and signed lateral edges (the $\sigma$ structure) as a model
        of excitatory/inhibitory lateral interaction.
\end{itemize}
The hypothesis is falsifiable and sharp: \emph{at matched parameters, the
G\"omb--Soma per-depth features predict V1/V2/V4 fMRI RDMs at least as well as a
ResNet-tiny}, with a paired bootstrap over the 92 conditions giving the
significance. A negative result here --- no better than ResNet --- would retire
the cortical-prototype framing; a positive result would be the model's first
genuine win, on the axis it was actually designed for.

\section{Implementation status and what is missing}
The scoring pipeline is built and tested in
\texttt{signedkan\_wip/src/cortical/}: \texttt{CorticalFeatureExtractor}
(per-depth retinotopic binning), \texttt{BrainScorer} (PLS\,$\to$\,ridge\,$\to$\,
$k$-fold CV\,$\to$\,split-half noise ceiling), a parameter-matched
\texttt{ResNetTinyCortical} baseline, and a shape-faithful \emph{synthetic}
Cichy-92 generator that drives an end-to-end smoke (92 images $\times$ 16
subjects in $\sim$22\,s/seed, 21 unit tests). The API is deliberately
dataset-agnostic: ingesting the real data is ``one new adapter function''.

The one missing piece is the \textbf{real Cichy-92 fMRI data}. The synthetic
generator validates shapes and the pipeline, but \emph{cannot} test the
hypothesis: whichever model scores higher on generated RDMs is an artefact of the
generator. Running the real benchmark requires (i) the published Cichy-92 fMRI
RDMs (or beta maps) in a known format, and (ii) the $\sim$50-line adapter that
maps them onto the existing \texttt{BrainScorer} interface. Both are small once
the data source is fixed.

\section{Conclusion}
G\"omb--Soma should not be defended as a classifier --- the evidence says it is
not one. Its honest evaluation is brain-predictivity, and Cichy-92 is the compact,
standard instrument for it. The protocol, baseline and code are ready; the next
concrete step is acquiring the real fMRI RDMs and writing the adapter, after
which a single parameter-matched run settles whether the cortical-prototype
framing survives.

\begin{thebibliography}{9}
\bibitem{cichy2014} R.~M. Cichy, D.~Pantazis, A.~Oliva.
  \emph{Resolving human object recognition in space and time.}
  Nature Neuroscience, 17(3):455--462, 2014.
\bibitem{kriegeskorte2008} N.~Kriegeskorte, M.~Mur, P.~Bandettini.
  \emph{Representational similarity analysis -- connecting the branches of
  systems neuroscience.} Frontiers in Systems Neuroscience, 2:4, 2008.
\bibitem{yamins2014} D.~L.~K. Yamins, H.~Hong, C.~Cadieu, E.~Solomon, D.~Seibert,
  J.~DiCarlo. \emph{Performance-optimized hierarchical models predict neural
  responses in higher visual cortex.} PNAS, 111(23):8619--8624, 2014.
\bibitem{schrimpf2018} M.~Schrimpf et al.
  \emph{Brain-Score: which artificial neural network for object recognition is
  most brain-like?} bioRxiv 407007, 2018 (and subsequent Neuron, 2020).
\bibitem{forman2003} R.~Forman.
  \emph{Bochner's method for cell complexes and combinatorial Ricci curvature.}
  Discrete \& Computational Geometry, 29(3):323--374, 2003.
\end{thebibliography}

\noindent\footnotesize\emph{Note on citations:} bibliographic details follow the
standard references for each method; verify exact pages/venues against the
project bib files (\texttt{paper/*/references.bib}) before external use.
\end{document}
